\chapter*{\chapterstyle{V --- Interpolation}}
\addcontentsline{toc}{section}{Interpolation}
Dans ce chapitre, on cherche à approximer par un polynôme une fonction donnée par un nuage de \(n + 1\) points de la forme \(\{(x_0, f(x_0)), \ldots,(x_n, f(x_n))\}\), on appelera une telle approximation \textbf{interpolation de la fonction} et on appelera un tel polynôme \textbf{polynôme interpolateur}.\<

On peut alors montrer qu'il existe \textbf{un unique tel polynôme} de degré \textbf{inférieur ou égal à \(n\)}, en effet, par 2 points passe une unique droite, par 3 points une unique parabole (potentiellement une droite si les points sont alignés) etc ...\<

L'objectif de ce chapitre sera donc de comprendre les différentes méthodes de construction d'un tel polynôme.

\subsection*{\subsecstyle{Système de Vandermonde {:}}}
On cherche donc un polynôme \(P\) de degré au plus \(n\) qui vérifie un système de \(n\) contraintes, en particulier pour \(P = \sum_{k=0}^{n} a_kX^k\), on a:
\[
   \begin{cases}
      a_0 + a_1x_0 + \ldots + a_nx_0^n = f(x_0)\\
      a_0 + a_1x_1 + \ldots + a_nx_1^n = f(x_1)\\
      \vdots\\
      a_0 + a_1x_{n-1} + \ldots + a_nx_{n-1}^n = f(x_{n-1})\\
      a_0 + a_1x_{n} + \ldots + a_nx_{n}^n = f(x_{n})
   \end{cases}
\]
Qui se ramène à l'équation matricielle:
\[
   \left(\begin{array}{cccc}
      1 & x_0 & \ldots & x_0^n\\ 
      1 & x_1 & \ldots & x_1^n\\
      \vdots & \vdots & \ddots & \vdots\\
      1 & x_{n-1} & \ldots & x_{n-1}^{n}\\ 
      1 & x_{n} & \ldots & x_{n}^n
   \end{array}\right) 
   \left(\begin{array}{c}
      a_0\\ 
      a_1\\
      \vdots\\
      a_{n-1}\\ 
      a_{n-1}
   \end{array}\right) =
   \left(\begin{array}{c}
      f(x_0)\\ 
      f(x_1)\\
      \vdots\\
      f(x_{n-1})\\ 
      f(x_{n})
   \end{array}\right)
\]
Qui se ramène donc à inverser une \textbf{matrice de Vandermonde} qu'on sait inversible, donc ceci permet de prouver l'existence et l'unicité d'un tel polynôme d'interpolation.

\begin{center}
   \textbf{Attention: Cette méthode est trés inefficace et trés couteuse en calculs !}
\end{center}
\subsection*{\subsecstyle{Interpolation de Lagrange {:}}}
Etant donné un nuage de \(n + 1\) abscisses \((x_0, \ldots, x_n)\), on définit alors le \(k\)-ième polynôme de la base de Lagrange (associé à l'abscisse \(x_k\)) par:
\customBox{width=5cm}{
   \[ L_{x_k} = \prod_{i \neq k} \frac{(X - x_i)}{x_k - x_i} \]
}
\begin{center}
   \textit{C'est un polynôme construit pour s'annuler en tout les points sauf en \(x_k\) et pour valoir \(1\) en \(x_k\).}
\end{center}
La construction d'un tel polynôme fait sens car alors, si on considère le polynôme:
\[
   P = \sum_{k=0}^{n}f(x_k)L_{x_k}
\]
On remarque que c'est bien un polynôme de degré au plus \(n\), et qu'on a bien par construction: 
\[
   \forall k \in \inticc{0}{n} , P(x_k) = f(x_k)
\] 
Donc \(P\) est bien le polynôme interpolateur de notre nuage de points.

\subsection*{\subsecstyle{Interpolation de Newton {:}}}
La dernière méthode d'interpolation est une méthode trés intéressante pour son affinité avec la programmation, en effet c'est une méthode \textbf{incrémentale}, c'est à dire que si on a un polynôme d'interpolation sur \(3\) points, c'est trés simple et peu couteux d'interpoler avec un nouveau point et d'obtenir donc un polynôme de degré supérieur.\<

Etant donné un nuage de \(n + 1\) abscisses \((x_0, \ldots, x_n)\), on définit alors le \(k\)-ième polynôme de la base de Newton par:
\customBox{width=5cm}{
   \[ N_k = \prod_{i=0}^{k - 1} (X - x_i) \]
}

Si on considère un nuage de points composé d'un seul point \((x_0, f(x_0))\), alors le polynôme d'interpolation est trivialement:
\[
   P_0 = f(x_0)
\]
Alors pour \textbf{pour rajouter un point} sans changer les valeurs précédentes\footnote[1]{En effet, la valeur en \(x_0\) ne change pas du fait du facteur \((x-x_0)\) et il suffit alors de résoudre pour \(\alpha_1\).
}, il suffit de poser:
\[
   P_1 = f(x_0) + \alpha_1(x - x_0)
\]
Alors pour \textbf{pour rajouter encore un point} sans changer les valeurs précédentes\footnote[2]{En effet, la valeur en \(x_0\) et en \(x_1\) ne change pas du fait du facteur \((x-x_0)(x-x_1)\) et il suffit alors de résoudre pour \(\alpha_2\).}, il suffit de poser:
\[
   P_2 = f(x_0) + \alpha_1(x - x_0) + \alpha_2(x-x_0)(x-x_1)
\]
On définit alors par récurrence le \textbf{polynôme d'interpolation} de degré \(n\) par:
\[
   P_n = P_{n-1} + \alpha_{n}N_n
\]

On appelle les coefficients \(\alpha_i\) les \textbf{différences divisées} d'ordre \(i\) de \(f\) et on les note:
\customBox{width=4cm}{
   \(\alpha_i = f[x_0, \ldots, x_i]\)
}
On peut alors montrer la \textbf{relation de récurrence} vérifiée par les différences divisées:
\[
   f[x_0] = f(x_0) \text{ et } f[x_0, \ldots, x_n] = \frac{f[x_1, \ldots, x_n] - f[x_0, \ldots, x_{n-1}]}{x_n-x_0}   
\]
Cette relation nous donne alors une nouvelle manière de construire le polynôme d'interpolation, en effet on a d'aprés l'expression par récurrence:
\[
   P_n = \sum_{k=0}^{n} f[x_0, \ldots, x_k] N_k   
\]
Il suffit donc de calculer les \(n\) différences divisées via la relation de récurrence et ainsi directement écrire la combinaison linéaire.\<

\underline{Exemple:} Si on considère le nuage de points \(\{(0, 2), (3, 7), (5, 1)\}\), alors on calculer les différences divisées dans un tableau:
\[
   \left(\begin{array}{cccc}
      x_0 & f[x_0] &  & \\ 
      x_1 & f[x_1] & f[x_0, x_1] &\\
      x_2 & f[x_2] & f[x_1, x_2] & f[x_0, x_1, x_2]\\ 
   \end{array}\right)  ; 
   \left(\begin{array}{cccc}
      0 & 2 &  & \\ 
      3 & 7 & \frac{5}{3} &\\
      5 & 1 & -3 & -\frac{14}{15}\\ 
   \end{array}\right)
\] 
Et donc \(P = 2 + \frac{5}{3}(X - x_0) -  \frac{7}{3}(X - x_0)(X - x_1) = 2 + \frac{5}{3}(X) -  \frac{14}{15}(X - 3)(X - 5)\)
\subsection*{\subsecstyle{Calculs de l'erreur {:}}}
On a la formule de majoration suivante pour l'erreur d'interpolation:
\customBox{width=12cm}{
   \[   
      |f(x) - P(x)| \leq \underset{m_x < \xi < M_x}{\sup}\Biggl\{\frac{f^{(n+1)}(t)}{(n+1)!}\Biggl\} |(x-x_0)(x-x_1)\ldots(x-x_n)|
   \]
}
Avec \(m_x = \min(x, x_0, \ldots, x_n)\) et \(M_x = \max(x, x_0, \ldots, x_n)\).\<

\underline{Exemple:} On cherche à interpoler la fonction logarithme avec le nuage de points \(\{(2, \ln(2)), (4, \ln(4)), (5, \ln(5))\}\). 
On calcule le polynôme d'interpolation et on obtient \(P(x) = \frac{-1}{5}()\)

\chapter*{\chapterstyle{V --- Dérivation Numérique}}
Dans ce chapitre, on cherche à approcher \textbf{les dérivées} d'une fonction \(f\) en un point \(x\) étant donné les images de la fonction en \(n + 1\) points \((x_0, \ldots, x_n)\).

\subsection*{\subsecstyle{Première approche {:}}}
Une première approche serait d'utiliser le \textbf{le polynôme interpolateur de Lagrange}, en effet, étant donnée un polynôme interpolateur \(P\), on peut montrer que:
\customBox{width=4cm}{
   \(f^{(k)}(x) \approx P^{(k)}(x)\)
}
\begin{center}
   \textit{Cette approche est conceptuellement trés simple mais couteuse en calculs et spécifique à la fonction étudiée.}
\end{center}

\subsection*{\subsecstyle{Seconde approche {:}}}
Le point clé de ce chapitre étant qu'une dérivée est liée \textbf{linéairement} aux valeurs de la fonction, et donc on cherche une expression telle que:
\customBox{width=16.6cm}{
   \begin{flalign*}
      f'(x) = \alpha_0f(x_0) + \alpha_1f(x_1) + \ldots + \alpha_nf(x_n)  \shorteqnote{(Pour certains coefficients \(\alpha_0, \ldots, \alpha_n\))}
   \end{flalign*}
   \vspace{-14pt}
}
Dans la suite par souci de concision, on travaillera sur un exemple précis, qui se généralisera facilement, on considère trois abcisses \(0, h, 2h\) et on cherche à approximer la dérivée de \(f\) en \(0\).\<

On peut alors calculer ces coefficients en utilisant \textbf{les formules de Taylor}, en effet on a:
\[
   \begin{cases}
      f(0) &= f(0) \\
      \color{BrightBlue1}f(h) &\color{BrightBlue1}= f(0) + hf'(0) + \frac{h^2}{2}f''(0) + o(h^2) \\
      \color{BrightRed1}f(2h) &\color{BrightRed1}= f(0) + 2hf'(0) + \frac{4h^2}{2}f''(0) + o(h^2)
   \end{cases}
\]
On remplace alors ces expressions dans la formule recherchée ce qui nous donne:
\begin{flalign*}
   f'(0) &= \alpha_0f(0) + \alpha_1{\color{BrightBlue1}f(h)} + \alpha_2{\color{BrightRed1}f(2h)}\\
         &= \alpha_0(f(0)) + \alpha_1{\color{BrightBlue1}(f(0) + hf'(0) + \frac{h^2}{2}f''(0))} + \alpha_2{\color{BrightRed1}(f(0) + 2hf'(0) + \frac{4h^2}{2}f''(0))} + o(h^2) \\
         &= f(0)(\alpha_0 + \alpha_1 + \alpha_2) + f'(0)(h\alpha_1 + 2h\alpha_2 ) + f''(0)\left(\frac{h^2}{2}\alpha_1  + 2h^2\alpha_2\right) + o(h^2)
\end{flalign*}
La dernière ligne nous donne alors un système à résoudre pour \(\alpha_0, \alpha_1, \alpha_2\) en fonction de \(h\), ie on veut:
\[
   \left(\begin{array}{ccc}
      \alpha_0 & \alpha_1 & \alpha_2\\ 
      0 & h\alpha_1 & 2h\alpha_2\\
      0 & \frac{h^2}{2}\alpha_1 & 2h^2\alpha_2
   \end{array}\right) 
   \left(\begin{array}{c}
      f(0)\\ 
      f'(0)\\
      f''(0)
   \end{array}\right) =
   \left(\begin{array}{c}
      0\\ 
      1\\
      0
   \end{array}\right)   
\]
Tout calculs faits, on obtient que \(f'(0) = -\frac{3}{2h}f(0) + \frac{2}{h}f(h) -\frac{1}{2h}f(2h) + o(h^2)\) qui est bien une approximaxion de la dérivée de \(f\) en \(0\).
\begin{center}
   \textit{Cette approche est conceptuellement assez complexe mais permet d'avoir une approximation générale d'une fonction quelconque.}
\end{center}
\pagebreak
\subsection*{\subsecstyle{Calculs d'erreur {:}}}
On reprends l'exemple ci-dessus et on exprime le reste d'ordre \(2\) sous sa forme de Taylor-Lagrange, on a:
\[
   \begin{cases}
      f(0) &= f(0) \\
      f(h) &= f(0) + hf'(0) + \frac{h^2}{2}f''(0) + \frac{h^3}{6}f'''(\xi_1) \\
      f(2h) &= f(0) + 2hf'(0) + \frac{4h^2}{2}f''(0) + \frac{4h^3}{3}f'''(\xi_2)
   \end{cases} \quad\quad\quad \text{(Pour certains \(\xi_1, \xi_2\) dans \(\icc{0}{h}, \icc{0}{2h}\))}
\]
Donc si on reprends la formule utilisée en notant\footnote[1]{Approximation en \(0\) de la première dérivée.} \(\Delta_0^{(1)}\) l'approximation obtenue, puis en substituant les \((\alpha_i)\) obtenus, on obtient:
\begin{flalign*}
   f'(0) &= \alpha_0f(0) + \alpha_1f(h) + \alpha_2f(2h)\\
         &= \Delta_0^{(1)} + \alpha_1\frac{h^3}{6}f'''(\xi_1) + \alpha_2\frac{4h^3}{3}f'''(\xi_2)\\
         &= \Delta_0^{(1)} + \frac{h^2}{3}f'''(\xi_1) - \frac{2h^2}{3}f'''(\xi_2)
\end{flalign*}
Donc pour \(I\) un intervalle qui contient tout les points \(0, h, 2h\) on obtient que:
\[
   \left|f'(0) - \Delta_0^{(1)}\right| = \left|\frac{h^2}{3}f'''(\xi_1) - \frac{2h^2}{3}f'''(\xi_2)\right| \leq \underset{t \in I}{\sup} \; \left|f'''(t)\right| \left|\frac{h^2}{3} - \frac{2h^2}{3}\right| = \underset{t \in I}{\sup} \; \left|f'''(t)\right| \left|\frac{-h^2}{3}\right|
\]
Finalement on a majoré notre erreur par:
\customBox{width=6cm}{
   \[   \left| f'(0) - \Delta_0^{(1)} \right|  \leq \underset{t \in I}{\sup} \; \left|f'''(t)\right| \left(\frac{h^2}{3}\right)
   \]
}
On peut alors établir une formule générale de l'erreur (en considérant alors une famille \((x_i)\) de \(n+1\) abscisses, pour lesquels on determine des \((\alpha_i)\)) et on a l'erreur d'approximation au point \(x\) est donnée par:


\chapter*{\chapterstyle{V --- Résolution approchée}}
\addcontentsline{toc}{section}{Résolution approchée}